% Part: sets-functions-relations
% Chapter: infinite
% Section: dedekind-induction

\documentclass[../../../include/open-logic-section]{subfiles}

\begin{document}

\olfileid{sfr}{infinite}{induction}
\olsection[Volledige inductie]{Dedekindalgebra's en volledige inductie}

Het cruciale punt is nu dat een Dedekindalgebra---zelfs \emph{elke}
Dedekindalgebra---als vervanging voor de natuurlijke getallen kan
dienen. Dit volgt onmiddellijk uit het volgende eenvoudige resultaat:

\begin{thm}[Volledige inductie]\ollabel{thm:dedinfiniteinduction}
Laat $N, s, o$ samen een Dedekindalgebra vormen. Dan geldt voor elke
verzameling $X$:
	\begin{center}
		als $o \in X$ en $(\forall {n} \in N \cap X){s}({n}) \in X$, {dan} $N \subseteq X$.
	\end{center}
\end{thm}

\begin{proof}
Volgens de definitie van een Dedekindalgebra is $N =
\closureofunder{s}{o}$. Als zowel ${o} \in X$ als $(\forall {n} \in
N)(n \in X \lif {s}({n}) \in X)$ geldt, dan behoort het gekozen
element ook tot de onderliggende verzameling; dus geldt ${o} \in N
\cap X$. Neem nu ${n} \in N \cap X$. De aanname geeft
${s}({n}) \in X$, terwijl
${s}({n}) \in N$ omdat $s \colon N \to N$. Dus is $N \cap X$
$s$-gesloten. Deze verzameling bevat het gekozen element en bevat
daarom de kleinste dergelijke verzameling. Bijgevolg is $N =
\closureofunder{s}{o} \subseteq X$.
\end{proof}

Inductie is kenmerkend voor de natuurlijke getallen. De kern is dus dat
we uit elke Dedekind-oneindige verzameling een Dedekindalgebra kunnen
vormen en die algebra als vervanging voor de natuurlijke getallen
kunnen gebruiken.

In \olref{thm:dedinfiniteinduction} is inductie weliswaar in
\emph{verzamelingstheoretische} termen geformuleerd, maar we kunnen
het principe eenvoudig in een vertrouwder vorm weergeven:

\begin{cor}\ollabel{natinductionschema}
Laat $N, s, o$ samen een Dedekindalgebra vormen. Dan geldt voor elke
formule $\phi(x)$, waarin parameters mogen voorkomen:
\begin{center}
	als $\phi(o)$ en $(\forall {n} \in N)(\phi(n)\lif
	\phi({s}({n})))$, {dan} $(\forall n \in N)\phi(n)$
\end{center}
\end{cor}

\begin{proof}
Neem $X = \Setabs{n \in N}{\phi(n)}$ en pas nu
\olref{thm:dedinfiniteinduction} toe.
\end{proof}

In dit resultaat spraken we van een formule ``met parameters''.
Daarmee bedoelen we ruwweg dat we voor willekeurige objecten $c_1,
\ldots, c_k$ met $\phi(x, c_1, \ldots, c_k)$ kunnen werken. Preciezer
kunnen we het resultaat zonder het woord ``parameters'' als volgt
formuleren. Voor elke formule $\phi(x, v_1, \ldots, v_k)$ waarvan alle
vrije variabelen zijn weergegeven, geldt:
	\begin{align*}
			\forall v_1 \ldots \forall v_k((&\phi(o, v_1,\ldots, v_k) \land {}\\
			&	(\forall x \in N)(\phi(x,v_1, \ldots, v_k) \lif \phi(s(x), v_1,\ldots, v_k))) \lif {}\\
			&\hspace{3em} (\forall x \in N)\phi(x, v_1,\ldots, v_k))
	\end{align*}
De uitdrukking ``met parameters'' maakt een tekst dus duidelijk veel
leesbaarder. (In \olref[sth][][]{part} zullen we deze manier van spreken
vaak gebruiken.)

Terug naar Dedekindalgebra's. In elke Dedekindalgebra kunnen we ook de
gebruikelijke rekenkundige functies voor optelling, vermenigvuldiging
en machtsverheffing definiëren. Dit is echter niet triviaal: daarvoor
is de techniek van de \emph{recursieve definitie} nodig. Die techniek
zullen we pas veel later invoeren en rechtvaardigen, in een veel
algemenere context. (Wie geïnteresseerd is, kan dit na
\olref[sth][ord-arithmetic][]{chap} opnieuw bekijken, of bijvoorbeeld
een alternatieve behandeling lezen, zoals \citealt[pp.~95--8]{Potter2004}.)
Als $N, s, o$ samen een Dedekindalgebra vormen, zullen we uiteindelijk
het volgende kunnen vastleggen:
\begin{align*}
	{a} + {o} &= {a} & & & {a} \times {o} &= {o} & & & {a}^{o} &= s(o)\\
{a} + {s}({b}) &= {s}({a}+{b}) &&& {a} \times {s}({b}) &= ({a}\times {b}) + {a}  & & & {a}^{{s}({b})} &= {a}^{b} \times {a}
\end{align*}
en aantonen dat deze functies zich gedragen zoals gewenst.

\end{document}
